% ----------------------------------------------------------
% Capitulo de revisão de literatura
% ----------------------------------------------------------
\chapter{Revisão bibliográfica}
\label{cap:revisao}
% ----------------------------------------------------------

A análise de redes complexas oferece perspectiva complementar para investigação de sistemas competitivos, permitindo modelar relações entre competidores através de grafos direcionados e ponderados. A abordagem fundamenta-se em arcabouço teórico consolidado que integra teoria de grafos, métricas de centralidade e algoritmos de detecção de comunidades. Este capítulo revisa os conceitos fundamentais e metodologias que embasam as escolhas técnicas adotadas neste trabalho.

\section{Fundamentação Teórica em Redes Complexas}

A teoria de redes complexas emergiu como um paradigma unificador para compreensão de sistemas intrincados em múltiplas disciplinas. A capacidade de modelar entidades como nós conectados por relações específicas oferece insights sobre propriedades emergentes que permanecem ocultas em análises tradicionais \citeonline{Ahmed2023ComplexNetworks}. Esta abordagem transcendeu fronteiras disciplinares, encontrando aplicação em sistemas biológicos, sociais, tecnológicos e econômicos. A robustez metodológica da análise de redes reside na flexibilidade de representação: diferentes definições de nós e arestas permitem capturar diversos tipos de interação dentro da mesma base conceitual. Esta versatilidade torna-se especialmente valiosa no contexto esportivo, onde competidores podem ser conectados através de confrontos diretos, proximidade de resultados ou relações de dominância competitiva.

As redes podem ser representadas matematicamente como grafos $G = (V, E)$, onde $V$ é o conjunto de vértices (nós) e $E$ é o conjunto de arestas que conectam pares de vértices. Em grafos direcionados, cada aresta possui uma orientação específica, representando relações assimétricas entre entidades. O peso associado a cada aresta quantifica a intensidade ou importância da conexão. Esta representação é fundamental para modelar relações competitivas: a direção da aresta indica resultado de confronto (do perdedor ao vencedor), enquanto o peso reflete magnitude da derrota ou diferença hierárquica entre posições no pódio \citeonline{Leicht2008CommunityStructure}. Grafos direcionados e ponderados capturam tanto informação ordinal (quem derrotou quem) quanto cardinal (com que intensidade), permitindo análises mais ricas que abordagens binárias não-ponderadas.

Redes complexas exibem propriedades emergentes não presentes em seus componentes individuais. A distribuição de grau heterogênea caracteriza-se por alguns nós com muitas conexões enquanto a maioria possui poucas, fenômeno observado em diversas redes naturais e artificiais. O coeficiente de agrupamento elevado reflete tendência de vizinhos de um nó estarem conectados entre si, formando triângulos e estruturas localmente densas. A existência de comunidades manifesta-se através de grupos densamente conectados internamente mas esparsamente conectados entre si \citeonline{Fortunato2010CommunityDetection}. Estas propriedades estruturais revelam padrões organizacionais fundamentais em sistemas esportivos: atletas centrais que enfrentam muitos competidores ao longo de suas carreiras, grupos de atletas que competem frequentemente entre si formando núcleos de rivalidade, e formação de agrupamentos baseados em era temporal, região geográfica ou especialização técnica. A análise destas propriedades permite identificar padrões que permaneceriam invisíveis em abordagens centradas exclusivamente em performance individual isolada.

\section{Modelagem de Redes Competitivas Esportivas}

A transformação de dados esportivos em redes complexas requer decisões metodológicas fundamentais sobre como representar atletas, competições e resultados através de nós, arestas e pesos. Estas escolhas determinam quais aspectos da estrutura competitiva serão capturados pela análise subsequente. A modelagem mais direta de competições esportivas define atletas como nós da rede. Esta representação é natural e intuitiva: cada participante das competições torna-se um vértice no grafo, preservando identidade individual ao longo de múltiplas edições. Alternativamente, algumas abordagens definem equipes como nós em esportes coletivos, ou agregam atletas por país para análises geopolíticas. A escolha de granularidade (individual vs coletiva) afeta fundamentalmente as propriedades estruturais da rede resultante e tipos de padrões detectáveis.

As arestas representam relações competitivas entre atletas. Múltiplas definições são possíveis: confrontos diretos em modalidades de eliminação, proximidade de resultados em provas cronometradas, ou compartilhamento de pódio em eventos com múltiplos medalhistas \citeonline{Radicchi2016OlympicDominance}. Este trabalho adota definição baseada em confrontos de pódio: uma aresta direcionada conecta dois atletas quando ambos conquistaram medalhas no mesmo evento, com direção apontando do atleta de posição inferior para o de posição superior. Esta escolha captura hierarquia competitiva direta: derrotar um campeão olímpico tem significado distinto de derrotar um competidor de nível inferior, informação preservada pela direção das arestas.

A quantificação de relações competitivas através de pesos numéricos constitui decisão metodológica crítica que afeta diretamente os resultados da análise de redes. A premissa fundamental é que nem todas as vitórias possuem mesmo valor competitivo: conquistar medalha de ouro representa realização qualitativamente superior a conquistar bronze, e esta diferença hierárquica deve ser capturada quantitativamente. A literatura em ranking olímpico estabelece sistema histórico de ponderação baseado na sequência de Fibonacci (3-2-1), onde medalha de ouro vale 3 pontos, prata vale 2 pontos, e bronze vale 1 ponto \citeonline{deSousa2012NetworkBased}. Este sistema, utilizado desde os Jogos Olímpicos de 1912 para classificação de países, fundamenta-se em propriedade matemática elegante: ouro (3) equivale à soma de prata e bronze (2+1), e ouro vale o dobro de bronze. Esta ponderação reflete hierarquia de valor absoluto das conquistas: cada medalha possui valor intrínseco independente do contexto competitivo.

Entretanto, a modelagem de redes competitivas demanda perspectiva distinta: não o valor absoluto de medalhas, mas a distância competitiva entre atletas que compartilham o pódio. A problemática central emerge: como quantificar a magnitude da diferença entre posições hierárquicas em um confronto direto? A hipótese é que a distância entre extremos (ouro-bronze) deve ser maior que entre posições adjacentes (prata-bronze ou ouro-prata), capturando intuição de que superar dois níveis hierárquicos representa gap competitivo superior a superar apenas um. Estendendo o princípio de Fibonacci para modelar distância ao invés de valor, este trabalho adota sistema 5-3-2 para construção de arestas: medalhista de ouro recebe aresta de peso 5 do medalhista de bronze, peso 3 da prata, enquanto medalhista de prata recebe peso 2 do bronze. Esta ponderação ampliada preserva as razões proporcionais do sistema original enquanto aumenta a diferenciação: a distância ouro-bronze (5) é 2.5 vezes maior que prata-bronze (2), refletindo gap hierárquico mais pronunciado. Matematicamente:

\[
\begin{cases}
w_{\mathrm{Bronze} \to \mathrm{Ouro}} = 5 \\
w_{\mathrm{Prata} \to \mathrm{Ouro}} = 3 \\
w_{\mathrm{Bronze} \to \mathrm{Prata}} = 2
\end{cases}
\]

Quando dois atletas compartilham pódio em múltiplos eventos ao longo da história olímpica, os pesos são acumulados ($W_{ij} = \sum_{k} w_k$), criando arestas com valores elevados que indicam rivalidade sistemática de longo prazo. Esta acumulação temporal levanta questão metodológica adicional: confrontos antigos devem receber mesmo peso que confrontos recentes? Motegi e Masuda argumentam que confrontos não devem receber pesos uniformes ao longo do tempo, propondo sistema de ranking dinâmico onde pontuações decaem exponencialmente \citeonline{deSousa2012NetworkBased}. A formulação matemática modifica o peso acumulado de arestas através de decay temporal:

\[
W_{ij} = \sum_{k} w_k(t_k) \cdot e^{-\lambda (T - t_k)}
\]

\noindent onde $W_{ij}$ é o peso final da aresta do atleta $i$ para $j$, $w_k(t_k)$ é o peso base do confronto $k$ ocorrido no tempo $t_k$, $T$ é o ano de referência (tipicamente o ano mais recente no dataset), e $\lambda$ é a taxa de decaimento anual (decay rate). Para $\lambda = 0$, o sistema reduz-se à soma simples ($W_{ij} = \sum_k w_k$); para $\lambda > 0$, confrontos antigos contribuem exponencialmente menos para o peso total. Em aplicação a tênis profissional com $\lambda = 0.05$ (decay de 5\% ao ano), o sistema dinâmico demonstrou acurácia preditiva superior (66\% versus 62-63\%), validando a relevância do tratamento temporal em torneios contínuos.

Entretanto, a necessidade de decay depende da estrutura subjacente dos dados. Quando a topologia da rede já captura implicitamente temporalidade --- como em competições quadrienais onde gerações de atletas raramente se sobrepõem --- sistemas estáticos podem ser suficientes. Esta observação sugere que validação de robustez deve ser específica ao contexto. A adoção do sistema 5-3-2 sem precedente direto na literatura, combinada com a escolha entre acumulação simples versus decay temporal, exige validação empírica rigorosa. Para tanto, Shiue et al. estabeleceram protocolo baseado em correlação de rankings através do coeficiente tau de Kendall ($\tau$) \citeonline{Shiue2025ApplyingPageRank}. A métrica tau de Kendall quantifica concordância entre duas ordenações através da fórmula:

\[
\tau = \frac{n_c - n_d}{\frac{1}{2}n(n-1)}
\]

\noindent onde $n_c$ é o número de pares concordantes (ordenados igualmente em ambos rankings), $n_d$ o número de pares discordantes, e $n$ o número total de elementos. Valores variam entre -1 (inversão completa) e +1 (concordância perfeita). O protocolo compara rankings gerados por diferentes parametrizações (e.g., múltiplas taxas $\lambda$, diferentes sistemas de pesos) e estabelece threshold de robustez: $\tau > 0.70$ indica concordância forte, validando que escolhas metodológicas não são arbitrárias mas produzem resultados consistentes. Esta validação quantitativa torna-se crítica para assegurar que as redes construídas capturam propriedades estruturais reais da competição olímpica, e não artefatos de parametrizações específicas. Os resultados desta validação para o sistema 5-3-2 adotado são apresentados no Capítulo~\ref{cap:resultados}.

\section{Métricas de Centralidade e Ranking em Redes}

A identificação de nós importantes em redes complexas constitui problema fundamental em análise estrutural. Diferentes métricas de centralidade capturam distintas noções de importância: conexão direta (grau), intermediação (betweenness), proximidade (closeness) e influência recursiva (PageRank). No contexto esportivo, estas métricas revelam atletas centrais sob perspectivas complementares. A centralidade de grau é a medida mais intuitiva de importância nodal: conta o número de conexões diretas de cada nó \citeonline{Freeman1979Centrality}. Em grafos direcionados, distinguem-se in-degree (arestas recebidas) e out-degree (arestas enviadas). No contexto de redes competitivas olímpicas, in-degree indica quantos competidores foram derrotados por um atleta, enquanto out-degree representa derrotas sofridas. Atletas com in-degree elevado dominaram muitos oponentes; aqueles com out-degree alto competiram frequentemente no pódio mas em posições inferiores. A razão in/out-degree oferece métrica sintética de performance relativa: valores próximos a zero indicam dominância (raramente derrotado), valores elevados indicam participação sem dominância (frequentemente superado).

Betweenness centrality quantifica quanto um nó atua como ponte entre outros pares de nós na rede \citeonline{Freeman1979Centrality}. Formalmente, para cada par de nós, calcula-se a fração de caminhos mínimos passando por um nó específico. Nós com betweenness elevado ocupam posições estruturalmente críticas: sua remoção fragmentaria significativamente a rede. Em redes esportivas, atletas com betweenness alto conectam diferentes grupos competitivos: podem representar pontes temporais (atletas com longevidade excepcional conectando gerações), pontes geográficas (atletas que competiram contra oponentes de múltiplas regiões), ou pontes técnicas (atletas versáteis competindo em múltiplas especializações).

PageRank representa avanço fundamental sobre métricas de grau ao incorporar importância recursiva: um nó torna-se importante não apenas por ter muitas conexões, mas por estar conectado a outros nós importantes \citeonline{Brin1998PageRank}. Desenvolvido originalmente para ranqueamento de páginas web, o algoritmo fundamenta-se na metáfora do random surfer proposta por \citeonline{Brin1998PageRank}: imagine um navegante que percorre a rede seguindo arestas aleatoriamente; o PageRank de um nó corresponde à probabilidade estacionária de encontrar o navegante naquele nó após muitas iterações. Seguindo a formulação original de \citeonline{Brin1998PageRank} adaptada para o contexto de redes direcionadas, o PageRank $PR(v)$ de um nó arbitrário $v$ é definido recursivamente por:

\[
PR(v) = \frac{1-d}{N} + d \sum_{u \in M(v)} \frac{PR(u)}{L(u)}
\]

onde $v$ representa o nó cuja importância está sendo calculada, $u$ representa cada nó predecessor (que possui aresta apontando para $v$), $M(v)$ é o conjunto completo de nós predecessores de $v$, $L(u)$ é o out-degree de cada nó $u$ (número de arestas saindo de $u$), $N$ é o número total de nós na rede, e $d$ é o damping factor (tipicamente 0.85) que representa probabilidade de continuar seguindo arestas versus saltar aleatoriamente para qualquer nó. Em redes esportivas, PageRank captura qualidade dos confrontos: derrotar um campeão olímpico múltiplo (que ele próprio derrotou muitos competidores fortes) confere PageRank mais elevado que derrotar um medalhista isolado. Esta propriedade recursiva produz rankings que diferem substancialmente de contagens simples de medalhas \citeonline{brasnam2024F1, Shiue2025ApplyingPageRank}. Atletas com múltiplas medalhas em eventos de baixa competitividade podem ter PageRank inferior a atletas com menos medalhas conquistadas contra oposição de elite. A métrica revela hierarquias de importância competitiva que métricas tradicionais baseadas em volume não capturam.

PageRank demonstrou acurácia preditiva superior em múltiplos contextos esportivos. A acurácia preditiva refere-se à capacidade do ranking em prever corretamente o vencedor de confrontos futuros: quando dois competidores se enfrentam, o sistema prevê vitória daquele com PageRank superior. Em tênis masculino profissional, sistemas de ranking dinâmicos baseados em PageRank alcançaram aproximadamente 66\% de predições corretas de vencedores de partidas versus 62-63\% de contrapartes estáticas \citeonline{deSousa2012NetworkBased}. Em beisebol colegial taiwanês, PageRank apresentou concordância forte (tau de Kendall > 0.70) com rankings oficiais e até 92.9\% de acurácia preditiva ao usar dados históricos de quatro temporadas anteriores para prever resultados subsequentes \citeonline{Shiue2025ApplyingPageRank}. A consistência temporal através de múltiplas temporadas valida PageRank como métrica robusta para avaliação de performance competitiva de longo prazo.

\section{Detecção e Caracterização de Comunidades}

Comunidades em redes complexas são definidas como grupos de nós densamente conectados internamente mas esparsamente conectados com o restante da rede \citeonline{Fortunato2010CommunityDetection}. A identificação destas estruturas modulares revela organização funcional subjacente: em redes sociais, comunidades correspondem a círculos sociais; em redes biológicas, a módulos funcionais; em redes esportivas, a agrupamentos baseados em era temporal, região geográfica ou estilo competitivo.

A qualidade de uma partição em comunidades é quantificada pela modularidade $Q$, que mede densidade de conexões dentro de comunidades comparada com densidade esperada em rede aleatória com mesma distribuição de grau. Formalmente:

\[
Q = \frac{1}{2m} \sum_{ij} \left[ A_{ij} - \frac{k_i k_j}{2m} \right] \delta(c_i, c_j)
\]

onde $A_{ij}$ é o peso da aresta entre nós $i$ e $j$, $k_i$ é o grau do nó $i$, $m$ é o número total de arestas, $c_i$ é a comunidade do nó $i$, e $\delta(c_i, c_j) = 1$ se $i$ e $j$ pertencem à mesma comunidade, 0 caso contrário. Valores de modularidade variam entre -0.5 e 1.0, com valores superiores a 0.3 tipicamente indicando estrutura comunitária significativa. Modularidade próxima a zero sugere ausência de estrutura modular ou rede extremamente densa onde fronteiras comunitárias são difusas.

A literatura apresenta duas abordagens principais para detecção de comunidades em redes direcionadas ponderadas. Algoritmos baseados em otimização de modularidade, como Louvain \citeonline{Blondel2008Louvain}, operam através de maximização gulosa hierárquica da função de modularidade. O método alterna entre duas fases: otimização local, onde cada nó é movido para a comunidade vizinha que maximiza ganho de modularidade $\Delta Q$, e agregação, onde comunidades tornam-se super-nós de meta-rede. Este processo hierárquico produz partições com diferentes resoluções até convergência, com eficiência computacional $O(n \log n)$ adequada para redes grandes. Embora desenvolvido para grafos não-direcionados, Louvain pode ser aplicado a redes direcionadas através de simetrização (transformar arestas direcionadas em bidirecionais) ou adaptação da fórmula de modularidade para incorporar assimetria \citeonline{Leicht2008CommunityStructure}.

Alternativamente, métodos baseados em teoria da informação preservam direcionalidade através de princípios distintos. O algoritmo Infomap \citeonline{Rosvall2008Infomap} utiliza equação de mapa para comprimir descrições de passeios aleatórios direcionados na rede, identificando comunidades que minimizam comprimento de descrição. Ao modelar fluxo de informação através de random walks que respeitam direção das arestas, Infomap potencialmente captura hierarquias de dominância mais fielmente que métodos baseados em simetrização. A escolha entre abordagens depende das propriedades estruturais da rede analisada e do fenômeno modelado: redes com fluxo direcional persistente podem beneficiar-se de preservação explícita de assimetria, enquanto redes onde direcionalidade é secundária podem ser adequadamente analisadas através de simetrização com ganhos em eficiência computacional e interpretabilidade de comunidades detectadas.

A detecção de comunidades em redes complexas identifica agrupamentos estruturais, mas a mera identificação de fronteiras comunitárias é insuficiente para compreensão profunda de suas propriedades distintivas. A caracterização multidimensional de comunidades requer métricas complementares que avaliem três dimensões fundamentais: perfil competitivo, conectividade estrutural e posicionamento hierárquico.

A avaliação de perfil competitivo em redes olímpicas demanda quantificação de dominância baseada em hierarquia de resultados. Conforme discutido na Seção 2.2, o sistema de ponderação Fibonacci (3-2-1) quantifica valor absoluto de medalhas para ranking de países. Este mesmo sistema pode ser aplicado para avaliar perfil competitivo de comunidades: cada comunidade recebe índice de dominância baseado na distribuição de medalhas de seus membros, transformando distribuição qualitativa (ouro/prata/bronze) em métrica numérica que reflete qualidade do portfolio competitivo do agrupamento. Este uso difere conceptualmente da ponderação de arestas (5-3-2): enquanto pesos de arestas modelam distância competitiva entre atletas individuais em confrontos diretos, o índice de dominância quantifica valor agregado das conquistas de uma comunidade inteira.

A qualidade de uma partição comunitária pode ser avaliada através do grau de segregação estrutural entre agrupamentos. Comunidades fortemente segregadas exibem alta densidade de conexões internas e baixa densidade de conexões externas, enquanto comunidades altamente conectadas entre si indicam estrutura menos modular \citeonline{Fortunato2010CommunityDetection}. Em redes direcionadas, como as competitivas olímpicas, a análise de conectividade inter-comunidade revela padrões de rivalidade estrutural: pares de comunidades com densidade excepcional de confrontos mútuos ao longo da história representam ``rivalidades sistemáticas'' que transcendem competições individuais isoladas \citeonline{Leicht2008CommunityStructure}. A razão entre conexões internas e externas quantifica este grau de segregação, permitindo identificar comunidades coesas versus altamente integradas.

A detecção de comunidades produz partição horizontal da rede, mas comunidades não são necessariamente equivalentes em importância estrutural. A hierarquização de comunidades baseada em métricas de centralidade agregadas revela estrutura vertical no conjunto de agrupamentos: comunidades ocupando posições centrais na topologia global (núcleos estruturais) versus comunidades periféricas marginalmente conectadas. Esta análise multi-escala demonstra que metadados agregados (como PageRank médio) correlacionam com importância estrutural, permitindo classificação hierárquica que independe de volume quantitativo \citeonline{Eriksson2022NetworkMetadata}. A hierarquia estrutural revela que centralidade não é propriedade exclusivamente individual: comunidades podem ocupar posições estratégicas de intermediação mesmo sem concentrar atletas de alto desempenho individual, identificando agrupamentos que funcionam como ``pontes'' entre eras ou regiões geográficas distintas.

\section{Seleção de Modalidades Esportivas}

A aplicação de análise de redes complexas ao contexto esportivo enfrenta desafio metodológico fundamental: como selecionar estrategicamente modalidades representativas para análise? O programa olímpico atual compreende mais de 60 modalidades distintas, cada uma com características estruturais, temporais e competitivas únicas. Esta diversidade, embora valiosa, torna impraticável a análise exaustiva de todas as modalidades em um único estudo. A literatura em análise de redes esportivas revela concentração desproporcional em modalidades populares como futebol, basquetebol e tênis, com limitada exploração de esportes individuais e suas diferentes características estruturais \citeonline{BibliometricOlympic2021}, comprometendo tanto generalização teórica quanto compreensão comparativa de diferentes tipos de interdependência competitiva \citeonline{TeamSportsNetwork2017}.

A fundamentação científica para seleção de modalidades deve considerar múltiplas dimensões relevantes para construção e análise de redes robustas. Três critérios principais orientam esta seleção: densidade competitiva (número de participantes afeta diretamente densidade e robustez das redes construídas \citeonline{Silva2019}), representatividade estrutural (necessidade de capturar diferentes tipos de interdependência competitiva \citeonline{Davids2013}), e objetividade de resultados (mensurabilidade precisa de performance garante robustez na construção de arestas ponderadas \citeonline{deSousa2012NetworkBased}).

A classificação de modalidades esportivas reconhece três tipos fundamentais de interdependência competitiva que produzem propriedades de rede distintivas \label{sec:taxonomia}. Esportes individuais caracterizam-se por competição direta onde performance depende exclusivamente de capacidades individuais, com interdependência manifestando-se apenas através de confrontos competitivos. Esportes coletivos apresentam alta interdependência interna e externa, com complexidade tática resultando em redes densas \citeonline{TeamSportsNetwork2017}. Esportes mistos combinam elementos individuais e coletivos, oferecendo perspectiva única para análise comparativa dentro da mesma modalidade. Esta taxonomia fornece base teórica robusta para seleção equilibrada que capture a diversidade essencial de fenômenos competitivos presentes no universo esportivo olímpico.

\section{Aplicação ao Contexto Olímpico}

Com base nos critérios estabelecidos e taxonomia proposta, a seleção de modalidades para análise de redes complexas deve priorizar representatividade entre as três categorias principais de interdependência competitiva. Modalidades puramente individuais com alta densidade de participação e mensurabilidade objetiva permitem análise de redes baseadas em proximidade de performance \citeonline{Silva2019}. Esportes coletivos genuínos com diferentes níveis de complexidade tática permitem análise comparativa de propriedades emergentes \citeonline{TeamSportsNetwork2017}. Modalidades que combinam provas individuais e coletivas oferecem oportunidade única para análise controlada, permitindo isolamento metodológico dos efeitos da interdependência na estrutura de rede. Esta abordagem resulta em portfolio equilibrado que maximiza representatividade teórica dentro de limitações práticas.

A adequação da seleção pode ser validada através das propriedades estruturais das redes resultantes. A literatura estabelece expectativas claras: redes de esportes individuais devem apresentar densidade menor com modularidade alta, refletindo competição localizada \citeonline{Silva2019}; redes de esportes coletivos devem exibir densidade maior com clustering elevado, consistente com interdependência entre equipes \citeonline{TeamSportsNetwork2017}; redes de esportes mistos devem permitir análise comparativa revelando como mudanças no formato competitivo afetam métricas estruturais. Esta diversidade de propriedades esperadas confirma que seleção baseada em critérios objetivos captura adequadamente a variedade de fenômenos competitivos presentes em competições olímpicas.

\section{Trabalhos Relacionados}

A aplicação de análise de redes complexas ao domínio esportivo tem crescido substancialmente na última década, com estudos explorando diferentes modalidades, métricas e contextos competitivos. Faria e Ferreira (2024) conduziram análise de redes sobre pilotos de Fórmula 1, construindo grafos baseados em resultados de corridas e aplicando PageRank \citeonline{brasnam2024F1}. O trabalho modela relações competitivas através de arestas direcionadas ponderadas, abordagem similar à adotada neste estudo, porém aplicada a contexto distinto: competições individuais em séries temporais contínuas versus eventos discretos quadrienais. A principal diferença metodológica reside na construção das redes: enquanto o estudo de F1 conecta pilotos baseado em desempenho relativo dentro de corridas, o presente trabalho estabelece conexões através de confrontos diretos no pódio olímpico, refletindo diferenças estruturais fundamentais entre calendário extenso com múltiplos eventos por temporada versus competições quadrienais. Outra distinção relevante concerne o escopo: Faria e Ferreira focam exclusivamente em ranking individual via PageRank, enquanto este estudo expande para incluir detecção de comunidades e caracterização multidimensional de agrupamentos estruturais.

Buldu et al. (2019) aplicaram modelagem de redes complexas e decomposição espectral para análise de performance em natação \citeonline{buldu2019complex}, construindo redes onde nadadores são conectados baseado em proximidade de tempos cronometrados. A abordagem difere substancialmente da metodologia adotada: as arestas baseiam-se em similaridade de performance ao invés de confrontos diretos no pódio, e a análise privilegia decomposição espectral ao invés de métricas de centralidade tradicionais. Estas escolhas refletem objetivos analíticos distintos: identificar agrupamentos baseados em níveis de performance absoluta versus focar em relações competitivas diretas e hierarquias de dominância. O presente trabalho complementa esta literatura ao oferecer perspectiva alternativa, analisando estrutura de rivalidades através de confrontos efetivos no pódio.

Vega-Oliveros et al. (2023) aplicaram detecção de comunidades a redes de jogadores de basquete NBA, identificando agrupamentos baseados em padrões de jogo e interações em quadra \citeonline{ComplexNetworksBasketball2023}. O presente trabalho distingue-se ao aplicar detecção de comunidades a nível inter-atletas através de múltiplas competições olímpicas, ao invés de intra-equipe em jogos individuais. Esta mudança de escala temporal e organizacional produz comunidades que refletem padrões competitivos de longo prazo abrangendo décadas e gerações de atletas. A caracterização multidimensional de comunidades emprega métricas consolidadas adaptadas ao contexto olímpico: entropia de Shannon \citeonline{Rajaram2017ShannonEntropy} para quantificar diversidade temporal e geográfica, coeficiente de Gini para medir concentração de medalhas, e coeficiente de variação de PageRank para avaliar heterogeneidade de performance dentro de cada comunidade detectada.

O presente estudo contribui para a literatura através de três aspectos principais: escopo multi-esporte permitindo comparações estruturais sistemáticas entre modalidades com diferentes tipos de interdependência competitiva; escala temporal ampliada abrangendo 125 anos de história olímpica, revelando dinâmicas de longo prazo; e aplicação integrada de múltiplas métricas estabelecidas para caracterização sistemática de comunidades, permitindo compreender suas características distintivas no contexto competitivo específico de cada modalidade.